\section{The \texttt{PSD\_Detector} McStas Component}
Position-sensitive gas-filled detector with gaseous thermal-neutron
converter (box, cylinder or 'banana').

\subsection*{Identification}
\begin{itemize}
  \item \textbf{Site:} 
  \item \textbf{Author:} Thorwald van Vuure
  \item \textbf{Origin:} ILL
  \item \textbf{Date:} Feb 21st, 2005
\end{itemize}

\subsection*{Description}
\begin{lstlisting}
An n times m pixel Position Sensitive Detector (PSD), box, cylinder or banana
filled with a mixture of thermal-neutron converter gas and stopping gas.
This model implements wires detection, including charge drift, parallax, etc.
The default is a simple 1D/2D position detector. However, setting
the parameter 'type' to "events" will save the signal as an event file which
contains all neutron parameters at detection points, including time. This is
especially suited for TOF detectors. Please use Histogrammer afterwards.
The gas creation event (neutron absorption in gas) is flagged as SCATTERED==1,
and the charge detection (signal) is flagged as SCATTERED==2.

GEOMETRY:
A flat rectangular box of given depth is simulated, or alternatively
a cylinder (when giving radius), or a horizontal cylindrical 'banana'
detector (when giving the width along the arc).
Box position is at its center, the 'banana' and cylinder have their
position at focus point, symetric in angular range for the 'banana'.

GAS SPECIFICATIONS:
The converter gas can be specified as He-3, BF3 or N2.
The filling pressure of converter gas must be specified along
with the zdepth: this gives an absorption probability (e.g. ~90% for
0.18 nm neutrons at zdepth 3 cm and He-3 pressure 5 bar).
The stopping gas can be specified as one of the following gases:
CF4, C3H8, CO2, Xe/TMA or He, or any for which a table of stopping
powers is available. NB: each stopping gas table has specific
data on a pair of particles (from thermal-neutron absorptions in
either He-3, BF3 or N2) in a specific gas. Unusual choices like N2
as a converter and He as a stopping gas probably require the
calculation of a new table.
Tip: to do a simulation in pure He-3, specify 0 bar for the stopping gas.
The pressures of the two gases together put a lower limit on the
achievable spatial resolution, no matter what pixel size is
specified.
Example: 1 bar CF4 gives ~2.5 mm spatial resolution, 3 bar CF4 gives
~1 mm. C3H8 requires more pressure for the same resolution, then
Xe/TMA, then CO2 and finally He.

IMPLEMENTED FEATURES:
Taken into account are the following effects:
- (rectangularly distributed) error in the recorded position due to
the range of the neutron capture products in the gas
- angular dependence of absorption efficiency
- border effect (events outside the sensitive volume tend to be
counted in the border pixels)
- parallax effect (with, if desired, an electrostatic lens to
partially correct for this - use "LensOn=1")
- wall effect
To correctly take the wall effect into account, a value must be
specified for the energy threshold for acceptance of an event as a
valid neutron event. This normally serves to discriminate from
gammas. Because these are not simulated in McStas, a value of 0 keV
can be chosen for this threshold if one is interested in seeing the
maximum number of neutrons; however, a more realistic value would be
100 keV. This implies seeing the maximum number of neutrons as well,
except in geometries with dominating wall effect.
The border effect is a considerably complex phenomenon: it depends
on the precise electric field configuration near the border. It is
simulated here simply by introducing a dead zone, gas-filled,
bordering the sensitive volume. Events in there can still cause
signals in the detector. The dead zone can cause a shadow in the
sensitive volume, just as in reality.
The algorithm simply adds all events on the first 'pixel' in the
dead zone to the border pixels. This is a coarse approximation,
because the physical effect depends on the orientation of the
tracks and the energy threshold setting. The bottom line is that
tracking the position of the Center Of Gravity of the charge
deposition in the detector does not allow for accurate modeling of
the border effect, and therefore the border pixels should be
ignored, just as in a real detector.
Note that specifying 0 width of the dead zone next to the
sensitive volume, apart from being unrealistic, does not allow
accurate simulation of the border effect: in this case, there will
be a relative lack of events on the border pixels due to the wall
effect.
This component determines the position of the Center Of Gravity of
the charge cloud in the detector. It does not simulate independent
channels.
Manufacturing imperfections such as wire diameter variations and
spread in electronic amplifier component properties are not simulated.
This should allow an experimenter to identify these manufacturing
imperfections in his physical machine and correct for them.

Example: PSD_Detector(xwidth=0.192, yheight=0.192, nx=64, ny=64,
zdepth=0.03, threshold=100, borderx=-1, bordery=-1,
PressureConv=5, PressureStop=1,
FN_Conv="Gas_tables/He3inHe.table", FN_Stop="Gas_tables/He3inCF4.table",
xChDivRelSigma=0, yChDivRelSigma=0,
filename="BIDIM19.psd")

Example: PSD_Detector(xwidth=0.26, yheight=0.26, nx=128, ny=128,
zdepth=0.03, threshold=100, borderx=-1, bordery=-1,
PressureConv=5, PressureStop=1,
FN_Conv="Gas_tables/He3inHe.table", FN_Stop="Gas_tables/He3inCF4.table",
xChDivRelSigma=0, yChDivRelSigma=0,
filename="BIDIM26.psd")

Example: PSD_Detector(radius=0.0127, yheight=0.20, nx=1, ny=128,
threshold=100, borderx=-1,
PressureConv=9.9, PressureStop=0.1,
FN_Conv="Gas_tables/He3inHe.table", FN_Stop="Gas_tables/He3inCO2.table",
yChDivRelSigma=0.006,
filename="ReuterStokes.psd")

Example: PSD_Detector(awidth=1.533, yheight=0.4, nx=640, ny=256,
radius=0.73, zdepth=0.032, dc=0.026, threshold=100,
borderx=-1, bordery=-1,
PressureConv=5, PressureStop=1,
FN_Conv="Gas_tables/He3inHe.table", FN_Stop="Gas_tables/He3inCF4.table",
xChDivRelSigma=0, yChDivRelSigma=0.0037,
LensOn=1, filename="D19.psd")
\end{lstlisting}

\subsection*{Input parameters}
Parameters in \textbf{boldface} are required; the others are optional.

\begin{longtable}{p{0.22\textwidth}p{0.12\textwidth}p{0.46\textwidth}p{0.14\textwidth}}
\toprule
\textbf{Name} & \textbf{Unit} & \textbf{Description} & \textbf{Default} \\
\midrule
\endhead
nx & 1 & Number of pixel columns & 128 \\
ny & 1 & Number of pixel rows & 128 \\
xwidth & m & Width of detector opening, in the case of a flat box & 0 \\
radius & m & Radius of detector for cylindrical/banana shape & 0 \\
awidth & m & Width of detector opening along the horizontal arc of the detector, measured along the inside arc of the sensitive volume in the case of a 'banana' detector & 0 \\
yheight & m & Height of detector opening & 0 \\
zdepth & m & Depth of the sensitive volume of the detector & 0.03 \\
threshold & keV & Energy threshold for acceptance of an event & 100 \\
PressureConv & bar & Pressure of gaseous thermal-neutron converter (e.g. He-3) & 5 \\
PressureStop & bar & Pressure of stopping gas & 1 \\
interpolate & 1 & Performs energy interpolation if true & 1 \\
p\_interact & 1 & Fraction of statistical events to be treated by component. Set it to 0 to use real absorption probability & 0 \\
verbose & 1 & Gives more information and spectrum & 0 \\
LensOn & 1 & set to 1 to simulate an electrostatic lens according to P. Van Esch, NIM A 540 (2005) pp. 361-367. & 1 \\
dc & m & Distance from the entrance window to the cathode plane, where the correction zone of the lens ends for banana shape and LensOn=1 & 0 \\
borderx & m & Width of (gas-filled) border on the side of the detector & -2 \\
bordery & m & Height of (gas-filled) border on the top/bottom of the detector giving rise to various border effects. Set to -2 to obtain the height of one pixel, -1 for the depth of the detector (default), or specify a non-negativedistance & -2 \\
xChDivRelSigma & 1 & Sigma of the error in position determination along the x axis (relative to xwidth), uniquely due to charge division readout. Set to 0 in case of independent channel readout along the x dimension & 0 \\
yChDivRelSigma & 1 & Relative sigma due to charge division in y direction, relative to yheight. & 0 \\
bufsize & 1 & Size of neutron output buffer default is 0, i.e. save all - recommended. & 0 \\
restore\_neutron & 1 & If set, the monitor does not influence the neutron state & 0 \\
angle & deg & Angular opening of detector in the case of a 'banana' detector, then overrides awidth parameter & 0 \\
type & string & If set to 'events' detector signal saves signal into an event file to be read e.g. by Histogrammer. This is to be used for TOF detectors. & 0 \\
filename & text & Name of file in which to store the detector image. The name of the component is used if file name is not specified & 0 \\
FN\_Conv & text & Filename of the stopping power table of the converter gas & "He3inHe.table" \\
FN\_Stop & text & Filename of the stopping power table of the stopping gas & "He3inCF4.table" \\
\bottomrule
\end{longtable}

\subsection*{Links}
\begin{itemize}
  \item \href{run:/home/willend/willend-McCode/mcstas-comps/contrib/PSD_Detector.comp}{Source code} for \texttt{PSD\_Detector.comp}.
  \item The test/example instrument \textless{}a href="../examples/Test\_PSD\_Detector.instr"\textgreater{}Test\_PSD\_Detector.instr\textless{}/a\textgreater{}.
\end{itemize}
\IfFileExists{PSD_Detector_static.tex}{\input{PSD_Detector_static.tex}}{}